\section{Fixed-budget review of handwritten digits}
\label{app:budgetref}

\paragraph{Cases and intervention.}
The task is to transcribe one handwritten digit. We use all 1,797 labeled
$8\times8$ images in scikit-learn's bundled \texttt{load\_digits} dataset.
The pixel arrays are distinct. A review consumes one slot, reveals the recorded
digit $y$, and replaces the proposed digit $a$ with $y$. An accepted case
retains $a$. We assume the reviewer returns the recorded digit without delay
or error. The feedback record keeps the original proposal and its corrected
label, so a reviewed mistake remains a failure of the original predictor.
We count both the slots consumed and the loss remaining after review. No
time or monetary value is assigned to a slot.

\paragraph{Frozen predictor and representation.}
A NumPy permutation with seed 4701 selects 397 training rows without reference
to the labels. We divide each pixel by 16 and fit multinomial logistic
regression with an intercept, $L_2$ penalty, $C=1$, the L-BFGS solver,
tolerance $10^{-8}$ and at most 500 iterations. The proposed digit is the
highest-probability class. For retrieval, we subtract the training-set mean
pixel vector and normalize each resulting vector to unit length. Every method
uses the same fixed proposals, class probabilities and representation. Training
rows serve only to fit the predictor and representation. The classifier made
92 errors on its 1,400 held-out images. All subsequent trials reuse this
model and pool.

\paragraph{Who sees which labels.}
For each seed from 0 through 11, a new label-independent permutation divides
the pool into four disjoint parts. The first 150 cases supply the same fully
reviewed initial memory to all methods. The next 150 are a fully labeled
calibration set, used only to set mixture weights and neighborhood sizes.
The following 500 cases arrive in ten batches of 50. Each method ranks a
complete batch, commits to five reviews, and receives only those five labels
before the next batch. Earlier arrivals break ties after scores are rounded
to 12 decimal places. Each adaptive history thus contains the same 150
initial labels and 50 additional labels selected by that method.

The final 600 cases arrive in twelve more batches. At this boundary we freeze
each history and evaluate two review policies on the same cases. Each policy
again receives five slots per batch. The evaluator uses the labels to count
losses and corrections. No future label updates a score, a weight or a memory.
The two future evaluations are counterfactual alternatives, each with 60
reviews. A replay allocates 410 post-training labels per policy (150 initial,
150 calibration, 50 stream and 60 future). The frozen classifier can operate
without the historical review labels. Equal access here isolates the choice
of cases to review at a given capacity.

\paragraph{Changing the cost of a mistake.}
During the stream every error costs one. We first retain this policy on the
future batches, then use the cost $|a-y|$ on those same batches. For example,
reading a 1 as a 9 now costs eight units, and reading an 8 as a 9 costs one.
The cost schedules are specified in advance and both can use the original
calibration labels. This models a change from class accuracy to accuracy of a
numerical entry. The distance cost is a chosen utility function. The observed
labels and classifier errors come from the dataset. A review changes the final
answer in the current case. Any benefit from acquiring feedback is measured
separately by scoring the same future batches with just the initial memory.

\paragraph{Comparators and calibration.}
Arrival order takes the first five cases in a randomly ordered batch.
Nearest failure ranks by maximum cosine similarity to a reviewed error.
The frozen rule uses $1-\max_c p(c\mid x)$ under unit loss. Under the
changed policy it uses $\sum_c p(c\mid x)|a-c|$, retaining the original
proposal $a$. Thus the classifier alone can respond to the cost change.

Two local votes use the same cosine neighbors and paid reviews. The outcome
vote averages the realized losses of the neighbors' own proposals. The
class-label vote estimates the digit distribution from their corrected labels
and applies the loss of the \emph{current} proposal to each of those labels.
For example, a correctly predicted neighboring 2 contributes zero past
error to the outcome vote, but supports review when the current proposal
is 7. A neighbor of cosine similarity $s$ has weight
$1/\max(1-s,10^{-8})$ before normalization. Each vote can combine its
estimate with the frozen score. For each objective separately, we choose
$k\in\{3,7,15\}$ and local weight $w\in\{0,0.25,0.5,0.75,1\}$ to minimize
mean squared error against realized loss on the 150 calibration cases,
using only the initial memory as neighbors. Ties favor smaller $k$, then
smaller $w$. These settings stay fixed for the rest of the replay.

\paragraph{Deployed BEM score and matched blends.}
We insert each reviewed embedding through
\texttt{BidirectionalExperienceMemory.add\_experience} and rank the batch by
its public \texttt{risk\_signal}. We use the package defaults for cosine
threshold (0.5) and risk sensitivity (0.8), with time decay disabled in
this undated replay. The risk score uses similar stored failures. Successes
are recorded but do not contribute to this read operation. The batch quota
is imposed by the experiment. We evaluate this risk component with a fixed
representation and without adapter updates.

For sign-only memory, success receives $+1$ and error receives $-1$.
For the changed cost, graded BEM re-encodes a reviewed error of distance $d$
as $-(0.35+0.65d/9)$. All nine possible distances pass the package's failure
threshold. Its severity buckets map distances 1--2, 3--6 and 7--9 to
severities 3, 4 and 5, respectively. We replay the exact same reviewed
records into the sign and graded memories. The encoding is computed entirely
from the paid historical labels.

We also give BEM the same opportunity to use the frozen score as either vote.
For both encodings we minimize the same calibration loss over the same five
weights in $wR+(1-w)F$, with $R$ the BEM score and $F$ the frozen score.
Under distance loss we first multiply $R$ by the mean distance of errors
in the initial memory, common to both encodings, to express the blend in
digit-distance units. The fallback scale for an error-free initial memory is
one. The sign and graded blends share the review history collected by the
unit-loss blend. Raw BEM collects its own history during the stream.
Comparisons between those two histories measure the complete review policy.

\paragraph{Joint use of failure, success and familiarity.}
A separate policy makes its own five stream reviews per batch. It calls the
same memory's \texttt{risk\_signal}, \texttt{success\_signal} and
\texttt{coverage\_signal} on each case when the calibrated joint weight is
positive. Write their outputs as $R,S,C$. Default coverage is the
package's calibrated cosine-kernel density score. We rank by
\[
  w h \{\alpha R+\beta(1-S)+\gamma(1-C)\}+(1-w)F .
\]
Here $h=1$ under unit loss and equals the initial mean distance per error
under distance loss. We use sign feedback for the unit objective and
re-encode the same policy history with graded outcomes for distance.
Calibration considers four ordered triples $(\alpha,\beta,\gamma)$, starting with
equal thirds, then $(1/2,1/4,1/4)$ and its two permutations emphasizing
lack of success or lack of coverage. Each component contributes when $w>0$.
We search the same five outer weights $w$ as the other blends, using all
20 combinations on the same calibration cases and initial memory. Ties
favor smaller $w$ and then the first triple in the order given. At $w=0$
the classifier alone ranks reviews, so the inner weights have no effect.
This is one specified way of combining the three package signals.

\begin{table}[htbp]
\centering
\caption{Paid reviews and their downstream outcomes. The stream columns cover
600 reviews and 412 underlying errors in 6,000 presentations. Each future
policy has 720 reviews, 454 underlying errors and 2,074 available
digit-distance units in 7,200 presentations. Counts pool 12 paired replays.}
\label{tab:digit-budget}
\small
\begin{tabular}{@{}lrrrrr@{}}
\toprule
 & Stream & \multicolumn{2}{c}{Future, unit cost} &
   \multicolumn{2}{c}{Future, distance cost} \\
\cmidrule(lr){3-4}\cmidrule(l){5-6}
Ranking & Caught & Caught & Left & Saved & Left \\
\midrule
Arrival order             &  33 &  41 & 413 &  184 & 1,890 \\
Nearest failure           & 143 & 176 & 278 &  944 & 1,130 \\
BEM risk, sign            & 140 & 178 & 276 &  959 & 1,115 \\
BEM risk, graded          & 140 & 178 & 276 &  975 & 1,099 \\
Frozen predictor          & 259 & 291 & 163 & 1,349 & 725 \\
BEM sign + frozen         & 251 & 291 & 163 & 1,454 & 620 \\
BEM graded + frozen       & 251 & 291 & 163 & 1,455 & 619 \\
Joint BEM + frozen        & 258 & 290 & 164 & 1,353 & 721 \\
Outcome vote + frozen     & 257 & 319 & 135 & 1,531 & 543 \\
Class-label vote + frozen & 303 & 379 &  75 & 1,693 & 381 \\
\bottomrule
\end{tabular}
\end{table}

\paragraph{Performance on future decisions.}
Table~\ref{tab:digit-budget} counts corrected errors under unit cost and
distance avoided under the changed cost. The class-label vote caught
379/454 errors (83.5\%), compared with 291/454 (64.1\%) for frozen confidence
and 178/454 (39.2\%) for raw BEM. Of the 6,480 automatically accepted
predictions, 75, 163 and 276, respectively, were wrong. With 720 perfect
reviews included, the corresponding end-to-end accuracies were 98.96\%,
97.74\% and 96.17\%. These totals count case presentations since each image
can recur in other seeds. The raw BEM advantage over nearest failure was
0.17 errors caught per seed (95\% interval $[-1.17,1.58]$).

Updating the class-label vote with its 50 paid stream labels changed 155 of
its 720 future review slots relative to its initial-memory version and
caught 41 more future errors, or 3.42 per seed $[1.83,5.00]$.
Raw BEM changed 200 slots and caught 25 more, or 2.08 per seed
$[0.83,3.33]$. In the matched BEM blend, 160 slots changed but the total
caught fell by three, a difference of $-0.25$ per seed $[-1.67,0.92]$.
Acquiring labels can change a review decision without improving its outcome.

\paragraph{What blending and grading changed.}
Under unit loss, calibration assigned positive weight to BEM in ten seeds
($w=0.25$ in seven and $w=0.5$ in three). The matched blend and frozen
confidence caught the same total, with a paired difference of 0.00 errors
per seed $[-2.25,2.08]$. The class-label and outcome votes caught 7.33
$[5.08,9.83]$ and 2.33 $[0.42,4.67]$ more errors per seed than the matched
BEM blend.

Under the changed cost, re-encoding raw BEM displaced 60 of 720 reviews
and saved 16 additional distance units, or 1.33 per seed $[-1.67,4.50]$.
The signed and graded BEM blends each used positive BEM weight in nine
seeds. The graded blend saved 106 more distance units than the frozen
expected-loss rule, or 8.83 per seed $[2.50,15.42]$, and one more unit
than the signed blend, or 0.08 per seed $[-2.83,3.50]$. The class-label
vote saved 19.83 more units per seed than the graded blend $[13.58,27.42]$.
The data support a gain over the frozen cost-sensitive rule for this blend
in these replays. The incremental contribution of graded encoding is
unresolved.

\paragraph{How often the joint score affected review.}
Calibration selected the classifier alone in nine of twelve unit-loss
replays. The joint weight was 0.25 twice and 0.5 once. Under distance loss
the joint weight was zero in eleven seeds and 0.75 in one. In every active
blend the inner weights were $(0.5,0.25,0.25)$ on failure risk, lack of
success and lack of coverage. At the same capacity, the policy caught
290 errors and avoided 1,353 distance units. Its paired differences from
the frozen rule were $-0.08$ errors per seed $[-0.50,0.25]$ and $+0.33$
distance units $[0.00,1.00]$. The risk-only graded blend saved 8.50 more
distance units per seed $[1.58,15.33]$ than this joint rule. These data
give little support for adding the three signals in this particular
combination to classifier confidence.

\begin{table}[htbp]
\centering\small
\caption{Paired future results by replay seed. Under unit loss the entries
count corrected errors out of the available errors. The final columns give
extra distance units saved by the graded BEM blend against each comparator
under the changed cost. Every cell uses the same budget of 60 reviews per method
and seed.}
\label{tab:digit-paired}
\begin{tabular}{@{}rrrrrrrr@{}}
\toprule
 & & \multicolumn{4}{c}{Errors caught} & \multicolumn{2}{c}{Graded blend gain} \\
\cmidrule(lr){3-6}\cmidrule(l){7-8}
Seed & Errors & Frozen & Risk blend & Joint blend & Class vote & vs.\ frozen & vs.\ sign blend \\
\midrule
0 & 40 & 28 & 23 & 28 & 30 & 0 & 0 \\
1 & 35 & 22 & 27 & 22 & 31 & 10 & 4 \\
2 & 42 & 20 & 23 & 20 & 36 & 26 & -6 \\
3 & 38 & 27 & 26 & 27 & 33 & 4 & -9 \\
4 & 36 & 21 & 26 & 21 & 32 & 30 & 0 \\
5 & 44 & 28 & 28 & 28 & 28 & 0 & 0 \\
6 & 32 & 24 & 20 & 22 & 28 & -13 & -5 \\
7 & 35 & 21 & 25 & 21 & 31 & 8 & 0 \\
8 & 35 & 23 & 24 & 23 & 29 & 13 & 3 \\
9 & 40 & 26 & 18 & 26 & 35 & 18 & 15 \\
10 & 31 & 21 & 21 & 21 & 26 & 0 & 0 \\
11 & 46 & 30 & 30 & 31 & 40 & 10 & -1 \\
\bottomrule
\end{tabular}
\end{table}

\paragraph{Uncertainty, scope and reproduction.}
Table~\ref{tab:digit-paired} shows how the effect varies by replay.
Each interval is the 2.5th and 97.5th percentile of 10,000 bootstrap means
of paired differences, resampling the twelve whole seeds with replacement
(NumPy seed 15001). All methods share the same batches within a draw.
The intervals describe sensitivity to random partitions and arrival order
on this finite corpus, conditional on the one trained classifier. They give
no independent replication across populations, writers or trained models.
The arrivals have simulated order, the reviewer is perfect, and the changed
cost is announced. There is no observed production shift, human review cost,
or evidence about language-model responses in this experiment.

From the release root with its environment active, reproduce the per-seed
counts, calibrated weights, selection comparisons and paired intervals with
\begin{verbatim}
python -m experiments.budget_benchmark --check results/budget.json
\end{verbatim}
The check uses the packaged digit data, one numerical thread and the
repository's BEM implementation. It requires exact counts and review-order
hashes, and tolerates floating-point roundoff in the reported summaries.
